%% hybrid.tex
%% Chapter: Hybrid Architectures for Unified Multimodal Large Language Models.
%% Designed to be \input{} into a master ACM acmart manuscript.
%% Bibliography file: hybrid.bib

\section{Hybrid Architectures}
\label{sec:hybrid}

\subsection{Motivation: Why Hybrid?}
\label{subsec:hybrid-motivation}

Unified multimodal large language models (MLLMs) face a representational dilemma. A purely discrete pipeline tokenizes images with a vector-quantized (VQ) codebook~\cite{DBLP:conf/nips/OordVK17,DBLP:conf/cvpr/EsserRO21,yu2024magvitv2} and treats them as a foreign language for the language model, as exemplified by Chameleon~\cite{Team202409818}. The discrete route enjoys a clean next-token training objective, yet it also discards fine-grained visual cues during quantization, which constrains comprehension performance below that of continuous CLIP-driven systems~\cite{Lin202505422,DBLP:conf/cvpr/YangYZRZ0Z25}. A purely continuous pipeline, in contrast, embeds images as dense features and trains a diffusion or flow head jointly with the language head, as in Transfusion~\cite{Zhou202411039} and Show-o~\cite{Xie202412528}. The continuous route preserves visual fidelity but couples generation tightly to a non-autoregressive backbone and complicates interleaved decoding.

Hybrid architectures emerge to reconcile these competing demands. They mix continuous and discrete signals along three axes: the tokenizer, the input--output pathway, and the network parameters that process each pathway. Section~\ref{subsec:hybrid-taxonomy} introduces a taxonomy that organizes recent hybrid designs along these axes; the remainder of this chapter then dissects each subclass and closes with a cross-cutting comparison and an outlook on open problems.

\subsection{A Taxonomy of Hybrid Designs}
\label{subsec:hybrid-taxonomy}

We organize hybrid architectures into five families, ordered from tokenizer-level mixtures to network-level mixtures. 

\begin{itemize}
    \item \textbf{Embedding Quantization} compresses continuous visual embeddings into a discrete code sequence aligned with text vocabulary, treating images as a foreign language~\cite{Ge202308041,Ge202308041,jin2024lavit,DBLP:conf/icml/Jin00XCJHSLZ0GM24,DBLP:conf/acl/ZhanDYZZLZYZL0F24,Ge202414396}.
    \item \textbf{Latent Semanticization} retains continuous tokens but injects high-level semantics through CLIP distillation, masked autoregression, or auxiliary objectives, narrowing the gap to discriminative encoders~\cite{DBLP:conf/cvpr/YangYZRZ0Z25,Lin202505422,Xie202515564,Huang202506590}.
    \item \textbf{Task-Specific Decoupling} splits understanding and generation across two encoders or two parameter branches inside the same transformer, following the Janus--JanusFlow--Janus-Pro lineage and its many descendants~\cite{DBLP:conf/cvpr/WuCWMLPLXYR025,DBLP:conf/cvpr/MaLCLWWPXZYZW0R25,Shi202415188,Chen202517811,Zhuang202512327,Fan202513436,Zhuang202512327,Tian202514682,Liao202505472,Deng202514683,Wu202518871,Wang202523044,Xu202510395,Lin202505422,Wang202523044,Wei202504548,Wang202523044,Tian202514682}.
    \item \textbf{Concatenation and Stacking} feeds heterogeneous tokens---continuous and discrete---into the language model along the sequence axis or the channel axis without sharing a single tokenizer~\cite{Huang202506590,Jiao202504423,He202504810}.
    \item \textbf{Fusion and Merging} unifies the two streams inside the tokenizer itself, producing a single index that simultaneously maps to a semantic feature and a pixel-level feature~\cite{DBLP:conf/iccv/XieDSL25,DBLP:conf/cvpr/QuZLWJGYDYW25,Chen202517811}.
\end{itemize}

The five families differ in \emph{where} the hybrid bridge sits. Embedding Quantization and Latent Semanticization fuse modalities at the tokenizer; Task-Specific Decoupling fuses at the network; Concatenation and Stacking fuses at the input sequence; and Fusion and Merging fuses inside the codebook. Each location implies a distinct trade-off between training simplicity, generation fidelity, and comprehension accuracy, which we examine in turn.

\subsection{Embedding Quantization}
\label{subsec:embed-quant}

Embedding-quantization methods train a tokenizer that compresses continuous visual features into a discrete code sequence consumable by a standard language-model loss. The seminal SEED tokenizer~\cite{Ge202308041} processes BLIP-2 ViT features through a causal Q-Former, applies vector quantization with a learned codebook, and supervises the codes with both image--text contrastive loss and image-reconstruction loss. The dual supervision lets the codes carry text-aligned semantics while still decoding into realistic pixels via a diffusion-based de-tokenizer. SEED-LLaMA~\cite{Ge202308041} couples this tokenizer with LLaMA-2 and trains the language model to predict the next image or text token, demonstrating that a single autoregressive objective suffices for both captioning and image synthesis.

Subsequent work scales this recipe in three directions. LaVIT~\cite{jin2024lavit} relaxes the fixed-length assumption with a dynamic visual tokenizer that adaptively merges redundant patches before quantization, producing variable-length code sequences whose entropy matches the underlying image complexity. Video-LaVIT~\cite{DBLP:conf/icml/Jin00XCJHSLZ0GM24} extends the design to video by decomposing each clip into a discrete keyframe code and a discrete motion code, which together capture the spatiotemporal dynamics that a frame-only tokenizer would lose. AnyGPT~\cite{DBLP:conf/acl/ZhanDYZZLZYZL0F24} pushes the language-as-foreign-modality view further, using independent discrete tokenizers for speech, music, and image, all of which feed a single LLaMA backbone. SEED-X~\cite{Ge202414396} returns to the image domain but enriches the de-tokenizer with multi-granularity conditioning, allowing the same discrete codes to drive both coarse generation and high-precision editing.

Across these works, the design knobs are the codebook size, the supervision target (CLIP alignment versus pixel reconstruction), and the quantization granularity (fixed grid versus dynamic merging). Larger codebooks improve reconstruction but slow convergence; CLIP alignment boosts comprehension but hurts pixel fidelity; dynamic merging improves efficiency but complicates the language modeling objective. The persistent bottleneck remains the quantization loss, which prompts the next family to keep tokens continuous while still injecting semantics.

\subsection{Latent Semanticization}
\label{subsec:latent-sem}

Latent-semanticization methods retain continuous visual tokens but reshape them so that they carry the high-level semantics that VQ-based pipelines lose during quantization. MMAR~\cite{DBLP:conf/cvpr/YangYZRZ0Z25} formalizes the rationale: it shows that diffusion-based unified models such as Transfusion encode image information in noise-conditional hidden states, which prevents the clean-input pass used during understanding from accessing the full visual content. MMAR therefore disentangles the diffusion head from the autoregressive backbone, attaches a lightweight diffusion head to each continuous patch embedding, and supervises the backbone hidden states with a v-prediction objective; this preserves all visual information for understanding while still supporting high-quality generation.

Other works restructure the tokenizer rather than the head. TokLIP~\cite{Lin202505422} reverses the usual VQ-then-CLIP pipeline by semanticizing existing VQ tokens: it re-encodes VQGAN codes with a ViT under text-alignment and CLIP-distillation losses, attaining LLaVA-comparable comprehension with less than 20\% of VILA-U's pretraining data. Show-o2~\cite{Xie202515564} builds unified visual representations on top of a 3D causal VAE through a dual-path of spatial--temporal fusion, which scales the original Show-o recipe to images and videos under a single language-and-flow head. Ming-UniVision~\cite{Huang202506590} pushes the agenda to its logical conclusion with a fully continuous tokenizer (MingTok) that progressively expands a compact latent into a high-dimensional semantic feature and finally a reconstructed pixel map, eliminating quantization error entirely.

These methods converge on a key insight: tokenizer-level fusion succeeds when the training objective explicitly disentangles ``what to generate'' from ``what to understand.'' Whereas Embedding Quantization compresses both signals into the same code, Latent Semanticization keeps the signals continuous and lets a downstream head decide which dimension to emphasize. Latent Semanticization therefore tends to outperform Embedding Quantization on comprehension benchmarks but pays the cost of more elaborate training pipelines and architecture-specific heads.

\subsection{Task-Specific Decoupling}
\label{subsec:decoupling}

Task-specific decoupling, the largest and most active hybrid family, addresses the optimization conflict between understanding and generation by allocating dedicated parameters to each task while sharing the language-model trunk. We organize this family by \emph{what} is decoupled: the visual encoder, the transformer branch, or the modeling paradigm.

\subsubsection{Dual Vision Encoders}
\label{subsubsec:dual-enc}

The Janus lineage decouples \emph{visual encoding}. Janus~\cite{DBLP:conf/cvpr/WuCWMLPLXYR025} uses SigLIP for understanding and a VQ tokenizer for generation, then routes both streams into a single autoregressive transformer; this resolves the granularity mismatch between the two tasks without duplicating the language model. JanusFlow~\cite{DBLP:conf/cvpr/MaLCLWWPXZYZW0R25} keeps the dual-encoder front end but replaces the discrete generation path with rectified flow, integrating diffusion-style sampling directly into the LLM and eliminating the need for a separate diffusion model. Janus-Pro~\cite{Chen202517811} scales the same blueprint with optimized training stages and 7B parameters, surpassing TokenFlow and MetaMorph on comprehension while approaching DALL-E 3 on instruction-following generation. The Skywork UniPic series~\cite{Wang202523044,Wei202504548}, Pisces~\cite{Xu202510395}, and Ovis-U1~\cite{Wang202523044} all inherit the dual-encoder backbone and add online reinforcement learning, mixed-resolution training, and bidirectional token refiners to push image-editing fidelity further.

\subsubsection{Dual Transformer Branches}
\label{subsubsec:dual-branch}

A second strand decouples the \emph{transformer parameters}. LMFusion~\cite{Shi202415188} retains a frozen Llama-3 for text and introduces parallel image-specific feed-forward and QKV projections; the two branches share only self-attention, so language ability is preserved while generation capability is added at half the FLOPs of training from scratch. BAGEL~\cite{Deng202514683} generalizes this idea to a Mixture-of-Transformer-Experts (MoT), in which an understanding expert processes ViT tokens and a generation expert processes VAE tokens, both routed through shared attention layers; this design powers a 7B-active-parameter model that exhibits emergent capabilities such as 3D navigation and free-form editing. Mogao~\cite{Liao202505472} adopts a closely related deep-fusion design with dual vision encoders and interleaved rotary position embeddings, while LightFusion~\cite{Wang202523044} demonstrates that the same architecture can be assembled by interleaving multimodal self-attention blocks between a frozen Qwen2.5-VL and a frozen Wan2.2 DiT, training only $\sim$35B tokens to reach competitive performance.

\subsubsection{AR plus Diffusion Hybrids}
\label{subsubsec:ar-diff}

A third strand decouples the \emph{generative paradigm}. VARGPT~\cite{Zhuang202512327} predicts text with next-token prediction and predicts images with next-scale prediction, fusing two autoregressive granularities inside one LLaVA-based backbone; VARGPT-v1.1~\cite{Zhuang202512327} adds DPO and instruction tuning to surface emergent editing behaviour. UniFluid~\cite{Fan202513436} keeps a pure autoregressive framework but uses continuous visual tokens with a diffusion head, reporting a clean trade-off between the two tasks that is controllable through a single loss-weight knob. UniGen~\cite{Tian202514682} and UniGen-1.5~\cite{Tian202514682} build on this paradigm with chain-of-thought verification at test time and unified RL rewards for generation and editing. OmniGen2~\cite{Wu202518871} and UniWorld-V1~\cite{Lin202505422} take a complementary route, attaching dedicated diffusion decoders to a multimodal understanding backbone via lightweight connectors and training only the generation pathway.

The decoupling family thus spans a spectrum from \emph{shallow} decoupling (different encoders, shared backbone) to \emph{deep} decoupling (different branches throughout the network, different generative objectives). Janus' empirical study established that even shallow decoupling delivers consistent gains, and BAGEL's MoT analysis confirms that deeper decoupling unlocks emergent abilities at scale. The cost is parameter and data efficiency: training BAGEL or Janus-Pro from scratch demands billions of multimodal tokens, which has motivated the lightweight assemblies of LightFusion and the data-efficient recipes of UniWorld-V1.

\subsection{Concatenation and Stacking}
\label{subsec:concat-stack}

Concatenation-and-stacking methods skip the tokenizer-level fusion and instead feed heterogeneous tokens into the language model in parallel. ILLUME+~\cite{Huang202506590} introduces DualViTok, a dual visual tokenizer that produces a semantic code and a pixel code per image; the language model consumes both codes in a coarse-to-fine sequence, generates discrete outputs, and offloads detail synthesis to a diffusion detokenizer. UniToken~\cite{Jiao202504423} concatenates VQGAN discrete tokens and SigLIP continuous embeddings along the sequence axis with explicit boundary tokens, explicitly preserving low-level details and high-level semantics for the same image. EMMA~\cite{He202504810} replaces sequence-axis concatenation with channel-axis concatenation; an aligned 32$\times$ autoencoder produces understanding and generation tokens of identical spatial shape, and the two streams fuse channel-wise inside a shared-and-decoupled transformer, which reduces context length while preserving cross-task interaction.

Concatenation-and-stacking succeeds because it is architecturally minimal: the language model sees a token sequence whose embeddings already carry mixed information, and no codebook redesign is needed. The trade-off is sequence length (in the sequence-axis variant) or training balance (in the channel-axis variant), both of which the next family addresses by collapsing the two streams into a single index.

\subsection{Fusion and Merging}
\label{subsec:fusion-merging}

Fusion-and-merging methods unify the two streams inside the tokenizer. MUSE-VL~\cite{DBLP:conf/iccv/XieDSL25} adds a semantic constraint to a VQGAN-style tokenizer through Semantic Discrete Encoding (SDE), so each visual code carries both pixel and semantic information; this yields a 4.8\% understanding gain over Emu3 with the same LLM size while cutting training data by roughly an order of magnitude. TokenFlow~\cite{DBLP:conf/cvpr/QuZLWJGYDYW25} adopts a dual-codebook architecture in which a semantic codebook and a pixel-level codebook share a common index, so a single token id retrieves two complementary embeddings; this separation enables the first discrete-input model to surpass LLaVA-1.5-13B on comprehension benchmarks. SemHiTok~\cite{Chen202517811} hierarchically organizes the codebooks: a semantic codebook is trained first and frozen, after which texture sub-codebooks are layered on top and supervise pixel reconstruction without disturbing semantic features.

Fusion-and-merging methods deliver the cleanest input format---one stream of discrete ids---at the price of a more involved tokenizer. Their success depends on disentangling the joint training of semantic and pixel objectives, which all three works address through dual codebooks, shared indices, or hierarchical supervision. The result is a tokenizer that approximates the expressivity of dual encoders while remaining compatible with off-the-shelf LLM training infrastructure.

\subsection{Cross-Cutting Comparison}
\label{subsec:hybrid-comparison}

Table~\ref{tab:hybrid-compare} summarizes the five hybrid families along four axes that recur throughout the literature: visual representation, training objective, decoupling depth, and headline trade-off. The comparison highlights three patterns. First, all five families converge on a continuous understanding stream and a discrete or continuous generation stream, but they differ in \emph{where} the streams meet---tokenizer, input sequence, or transformer branch. Second, deeper decoupling tends to yield higher peak performance at the cost of training data and parameters; shallow tokenizer-level fusion is more data efficient but caps the achievable comprehension. Third, the most recent works (BAGEL, Skywork UniPic 2.0, EMMA, LightFusion) combine multiple families: BAGEL pairs dual encoders with dual transformer branches; LightFusion pairs dual transformer branches with shallow attention fusion; EMMA pairs channel-wise concatenation with a shared-and-decoupled network. The hybrid landscape therefore moves from disjoint design choices toward composable design patterns.

\begin{table}[t]
  \caption{Cross-family comparison of hybrid architectures. ``D'' denotes discrete tokens and ``C'' denotes continuous tokens. Bridge column marks where the hybrid bridge sits: tokenizer (T), input sequence (S), or network branch (N).}
  \label{tab:hybrid-compare}
  \small
  \setlength{\tabcolsep}{4pt}
  \begin{tabular}{@{}p{0.18\linewidth}p{0.20\linewidth}p{0.18\linewidth}cp{0.24\linewidth}@{}}
    \toprule
    Family & Visual repr. & Training obj. & Bridge & Trade-off \\
    \midrule
    Embedding Quantization     & D (CLIP-aligned)            & next-token CE         & T & Codebook caps comprehension \\
    Latent Semanticization     & C with semantic supervision & CE + diffusion/flow   & T & Heavier tokenizer or aux. head \\
    Task-Specific Decoupling   & C (und.) + D/C (gen.)       & CE + diffusion/flow   & N & High peak, costly training \\
    Concatenation \& Stacking  & C $\oplus$ D                & CE + diffusion        & S & Long context or balance \\
    Fusion \& Merging          & shared index, dual codes    & joint reconstruction  & T & Intricate tokenizer training \\
    \bottomrule
  \end{tabular}
\end{table}

\subsection{Open Problems and Outlook}
\label{subsec:hybrid-outlook}

Hybrid architectures have closed much of the gap between unified MLLMs and specialist models, yet several questions remain open. Three directions stand out for future work.

\textbf{Tokenizer co-design with the LLM.} Most hybrid tokenizers are trained in isolation and then frozen during LLM training, which leaves the codebook misaligned with the evolving language priors. Ming-UniVision and SemHiTok point toward end-to-end co-training, but the optimization landscape of joint tokenizer--LLM training remains poorly understood.

\textbf{Scaling with mixed objectives.} BAGEL demonstrates that complex reasoning emerges only at sufficient scale and with sufficient interleaved data, yet the scaling laws for joint cross-entropy and diffusion losses are far from settled. UniFluid and EMMA report sharp loss-weight sensitivity, which suggests that adaptive curricula and modality-aware learning rates deserve systematic study.

\textbf{Modular composition of pretrained experts.} LightFusion shows that a competitive unified model can be assembled by interleaving attention blocks between a frozen VLM and a frozen DiT. If this trend continues, the next generation of hybrid models may resemble plug-and-play assemblies rather than monolithic systems trained from scratch, which would democratize access at the cost of new architectural-search problems.

The hybrid agenda thus reframes the unification of multimodal understanding and generation as a design-pattern problem rather than a single-architecture problem. As tokenizers, branches, and training objectives become composable, hybrid architectures are positioned to remain the dominant paradigm for unified MLLMs.

%% End of \section{Hybrid Architectures}
